\scanpage{15}
1) Morphisme d'un espace loc annelé dans schéma affine

\[
 \operatorname{Hom}_k(X,\underset{\operatorname{Spec} A}{\underset{\shortparallel}{Y}})
 \mathrel{\substack{\overset{\sim}{\longrightarrow}\\[-2pt]\longleftarrow}}
 \operatorname{Hom}_{k\text{-alg}}(\underset{\Gamma(Y,\underline O_Y)}{\underset{\shortparallel}{A}},\Gamma(X,\underline O_{\deleted{\unsure{x}}X}))
\]

\deleted{\(\forall x\in\)} Soit $u:A\mathrel{-}\Gamma(X,\underline O_X)$

Pour tout $x\in X$, soit $\varphi(x)$ \inserted{$\in Y=\operatorname{Spec}(A)$} le noyau de

\diagram{\begin{tikzcd}[ampersand replacement=\&]
 A\arrow[r,dash,"u"]\&\Gamma(X,\underline O_{\unsure{X|x|\mathcal X}})\arrow[rr,dash,"\varepsilon_x"]\arrow[dr]\&{}\&k(x)\\
 {}\&{}\&\underline O_{X,x}\arrow[ur]\&{}
\end{tikzcd}}

d'où $\varphi:X\longrightarrow Y$

Si $f\in A$, $\varphi^{-1}(U_f)={}$ ens des $x\in X$ tels que \inserted{$\varepsilon_x$}$(u(f))\ne 0$ i.e. tels que $u(f)_x$ inversible, i.e. $X_{u(f)}$ : il est ouvert donc $\varphi$ continu.

De plus, on trouve

\[ A_{\deleted{\unsure{g}}f}\longrightarrow\Gamma(U_f,\underline O_X)\qquad\text{rend $f$ inversible, \unsure{donc|donc se|en} \unsure{se factorise en|factorise en}} \]

\[ \Gamma(Y_f,\underline O_Y)\simeq A_f\mathrel{-}\Gamma(U_f,\underline O_X)\quad\text{avec compatibilité pour \unsure{restriction|restrictions}, d'où un morphisme annelé.} \]

Il est \uline{loc}. \uline{annelé}, car \deleted{\unsure{par définition, il transforme}} \deleted{\unsure{un}} on a $\varphi^{\unsure{-1|*}}(Y_f)=X_{u(f)}$ et de même après restriction à un ouvert \ldots

\uline{Conséquence} La catégorie $(\mathrm{Sch})_{/k}$ se réconstitue connaissant $\mathrm{Sch}=\mathrm{Sch}_{/\unsure{Z|\mathbb Z}}$ - et $S=\operatorname{Spec} k$ comme $(\mathrm{Sch})_{/S}$ : les $k$-schémas sont les $S$-schémas. Mais $\mathrm{Sch}_{/S}$ important sans que $S$ soit néc. affine.

\diagram{\begin{tikzcd}[ampersand replacement=\&,column sep=small,row sep=small]
 (\mathrm{Sch}\text{-Aff})_k^{\unsure{o|\circ}}\arrow[d,phantom,"\shortparallel"]\&{}\&{}\\
 (\mathrm{Aff}_k)\arrow[r,hook,"\mathrm{incl}"]\arrow[d,phantom,"\shortparallel"]\&(\mathrm{Sch})_{/k}\arrow[r,dash,hook']\&(\mathrm{Locan})_k\arrow[dl]\\
 (\mathrm{Alg}_k)^o\&\widetilde{(\mathrm{Sch})_{/k}}\arrow[r,hook,"\mathrm{restr.}","{\substack{\text{faisceaux}\\\text{pour ``top.}\\\text{de Zariski''}}}"']\&\widetilde{\mathrm{Aff}}_k\simeq\deleted{\unsure{\widetilde{\mathrm{Alg}}_k}}\underline{\operatorname{Hom}}^{\sim}(\mathrm{Alg}_k,\mathrm{Ens})\\
 \widetilde{(\mathrm{Aff}_k)}\arrow[uu,hook]\arrow[d,phantom,"\shortparallel"]\&{}\&{}\\
 \underline{\operatorname{Hom}}^{\sim}(\mathrm{Alg}_k,\mathrm{Ens})\&\widehat{(\mathrm{Sch})_{/k}}\arrow[uu,hook,"\mathrm{inc}"]\arrow[r,"\mathrm{restr.}","(\text{\unsure{pas pl. fid|pas pl. fid.|pas plein fid.}})"']\&\widehat{\mathrm{Aff}}_k\simeq\deleted{\unsure{\mathrm{Alg}|\mathrm{Alg}_k}}\underline{\operatorname{Hom}}(\mathrm{Alg}_k,\mathrm{Ens})\arrow[uu,hook,"\mathrm{inc}"]\\
 \widehat{(\mathrm{Aff}_k)}\arrow[uu,hook,"\mathrm{inc}."]\arrow[d,phantom,"\shortparallel"]\&{}\&{}\\
 \underline{\operatorname{Hom}}(\mathrm{Alg}_k,\mathrm{Ens})\&{}\&{}
\end{tikzcd}}

Comment exprimer en termes d'algèbre commutative qu'un \inserted{co-}foncteur sur $\mathrm{Aff}_k$ est un faisceau, ? i.e. un foncteur Aff

[ Interprétation en termes de foncteurs d'anneaux de

\diagram{\begin{tikzcd}[ampersand replacement=\&]
 {}\&(\mathrm{Sch})_k\&{\simeq}\&\mathrm{Sch}_{/\operatorname{Spec} k=S}\\
 C=\mathrm{Aff}_k\&\widehat{C_{/S}}\arrow[u,hook]\&{\simeq}\&\widehat C_{/\underline S}\arrow[u,hook]
\end{tikzcd}}
]

$\left\{\begin{aligned}&\text{Interprétation des \uline{sommes} en termes faisceautiques}\\
 &\text{(pas d'ennui pour les \deleted{produits} $\varprojlim$ (finies, \inserted{disons} \ldots.)}
 \end{aligned}\right.$

\margin{Notion de Sup de sous-faisceaux}

% archive sheet number: 14
